\section{Forces and Gravity (P8, P16, P21, P24, P28, P32, P33)}
\label{sec:gravity}

\subsection{The gauge force (P8)}

A force is the twist a tone gets crossing a note, a gauge connection on the links.
\textbf{Where we landed:} 3D SU(2) lattice gauge theory confines (the string tension is
positive at every coupling, falling from $1.32$ to $0.40$, the Wilson area law). The
chiral overlap fermion sees gauge topology exactly: in a U(1) background of charge $Q$
its index equals $-Q$, the lattice Atiyah-Singer theorem. And in a dynamical gauge
field a chiral condensate forms from the anomaly, in both the Abelian Schwinger model
and the non-Abelian SU(2) case. The open rung is the Weyl-projected chiral gauge theory,
unsolved in physics generally.

\subsection{The Newtonian limit (P16)}

The static potential is the Green's function of the emergent Laplacian. \textbf{Where we
landed:} it matches the continuum in every dimension, confining in 1D, logarithmic in
2D, and Newtonian $1/r$ in 3D (the $1/r$ form is the best fit at $R^2 = 0.997$, beating
$1/r^2$ and $\log$).

\subsection{The graviton, and its operator from the action (P21, P24)}

The graviton is the propagating excitation of the geometry, a symmetric rank-2 tensor.
\textbf{Where we landed:} it is a proper massless spin-2 field with exactly two
transverse-traceless polarizations (a massive spin-2 would have five), and a massless
dispersion. And its kinetic operator is derived from the action: the linearized Einstein
operator (the second variation of the Einstein-Hilbert action, the continuum limit of
the discrete Benincasa-Dowker action) is diffeomorphism-invariant (pure-gauge
perturbations are annihilated to machine zero) with exactly two massless modes at
eigenvalue $\tfrac12|k|^2$, no projector imposed.

\subsection{The Einstein equation as an equation of motion (P32)}

\textbf{Where we landed:} the linearized Einstein equation behaves as a real theory of
gravity. The Einstein tensor is transverse, $k \cdot G[h] = 0$ to machine zero, so the
equation automatically enforces energy-momentum conservation. It reduces to Newton in
the static limit (the Poisson equation, P16), and in vacuum it propagates a massless
graviton at the speed of light ($\omega/|k| = 1$ at every frequency). The full nonlinear
Einstein equation from the discrete action is the remaining step.

\subsection{Singularity resolution (P28)}

\textbf{Where we landed:} discreteness gives a minimum causal length (it scales as
$\rho^{-1/2}$ in 2D), so curvature, which goes as one over length squared, is capped at
a finite value that rises with density and never diverges. The big-bang and black-hole
singularities of general relativity are resolved: a finite-density beginning at the
discreteness scale, not an infinite-density point, which fits the bootstrap (the
universe starts from a first element, not a singular $t=0$).

\subsection{Black-hole entropy (P33)}

\textbf{Where we landed:} the entanglement entropy of a region scales with its horizon
\emph{area}, not its volume (a 3D area fit beats a volume fit). Reading the region as a
black hole and its boundary as the horizon, this is the Bekenstein-Hawking law
$S = A/4$, with the entanglement across the horizon as its microscopic origin. The full
black-hole physics (Hawking radiation, the exact coefficient, the information question)
is the long road ahead.
